\begin{abstract}
Can interpretable macroscopic PDEs be recovered from noisy
interacting-agent trajectories?
This thesis investigates the question for Vicsek--Morse particle systems
by pairing shift-aligned \gls{pod} forecasting (\gls{mvar} and LSTM
latent time-steppers) with weak-form PDE discovery (\gls{wsindy}) on
density and polarisation flux fields across nine dynamical regimes and
four initial-condition families.

WSINDy recovers sparse, regime-consistent effective PDEs---matching
the mechanism structure anticipated from kinetic theory
(advection, diffusion, Morse forcing)
with density weak-form $R^2>0.99$ in eight of nine regimes
(the ninth at $0.95$)---yet the
identified coefficients often reside outside the regime of numerical
stability for autonomous forward-integration.
This \emph{closure gap} between identification and rollout quality is
itself a diagnostic: the recovered coefficients are effective finite-$N$
parameters at the chosen KDE bandwidth, and the gap reveals which closure
ingredients the effective description still lacks.
In particular, the three-field library $(\rho,p_x,p_y)$ appears
insufficient to close the mean-field description when agent speed is
state-dependent, suggesting that variable-speed regimes require
additional closure fields beyond density and polarisation flux.

On the forecasting side, the $116\!\times$ parameter-count gap between
MVAR ($1{,}824$ parameters) and LSTM ($211{,}091$) does not yield
proportional accuracy gains: MVAR achieves forecast $R^2>0.99$ on the
strongest constant-speed regimes (CS\,+\,PV median~$0.995$), while
variable-speed variants degrade sharply
(VS median~$0.425$ for MVAR, $-0.052$ for LSTM), exposing a
shift-predictability bottleneck that neither model class overcomes.
Together, these findings delineate the boundary between what is
identifiable and what is predictable from coarse-grained
collective-motion data.%
\footnote{Code and configurations:
\url{https://github.com/eugenio-macias/wsindy-manifold}.}
\end{abstract}

\vspace{1.5em}
\noindent\textbf{Acknowledgments.}\enspace
I am grateful to my thesis advisor, Prof.\ Bj\"orn Sandstede, who gave me
the initial ideas that set the direction of this work and remained a source
of guidance as those ideas evolved into a full pipeline.
I am also grateful to my second reader, Prof.\ Ian~Wong, for his thoughtful
reading and encouragement.

This thesis was an opportunity to see the full power of applied mathematics
and modern computation working together.
Looking back, it draws on nearly every area I studied at Brown in the
APMA--CS--MATH track, often in ways I did not anticipate when I first
encountered those subjects.
Probability theory and statistics underlie the KDE fields that convert
discrete agent trajectories into continuous density snapshots, the noise
models, and the convergence analysis.
Linear algebra and functional analysis provide the backbone: POD and SVD
truncation, the Hilbert-space formulation of the weak form, and the
Kolmogorov $n$-width theory that explains \emph{why} transport-dominated
systems resist linear reduction.
Differential equations are everywhere --- in the ETDRK4 integrator, in
the PDE candidate library, in the Fokker--Planck framing of the mean-field
limit.
Machine learning and deep learning enter through the LSTM sequence models
and the SR3 sparse regression framework.
Abstract algebra makes a perhaps unexpected appearance: the continuous
SO(2) rotational symmetry of the system is precisely the kind of group
action studied in algebra, and the shift-alignment step is, in essence,
quotienting the data by that group action to obtain a canonical
representative.
Topology threads through the persistent homology literature against which
we position our approach, and statistics returns at the end in the form of
stability selection and model-order criteria.
I am grateful to all the faculty and courses at Brown that built these
foundations in me, and I find it exciting that a single undergraduate
thesis could call on all of them at once.
Applied mathematics is at its best when it refuses to stay in one lane ---
and high-performance computing gave us the experimental scale to take those
mathematical ideas seriously.

I thank the Center for Computation and Visualization
(CCV) at Brown University for providing access to
the OSCAR high-performance computing cluster, on which all particle
simulations, ROM training, WSINDy fitting, and ETDRK4 rollout computations
were performed.
